\subsection*{The nonzero time slope of a zero-left spectral Kummer correction}

\noindent
The frozen spectral repair preserves both left \emph{spatial} Cauchy
coordinates. Differentiating those identities at a moving left
endpoint does not preserve the zero \emph{temporal} Cauchy slope.
This distinction supplies the precise missing time interface for the
strictly verified weighted Volterra estimates.

Let \(\mathcal L_s=\partial_{rr}+p(s,r)\partial_r+q(s,r)\),
let \(\rho=\rho(s)>0\), and assume sufficient classical mixed
regularity on a swept positive-radius interval. For the same actual
source \(F\) and spectral displacement \(\delta\), set
\[
\mathcal L_su_0=F,\quad
u_0(s,\rho)=U,\quad (u_0)_r(s,\rho)=D,
\quad
\mathcal L_su_1=\delta u_0,\quad
u_1(s,\rho)=(u_1)_r(s,\rho)=0,\quad w=u_0-u_1.
\]
These are the defining differential equations of the original
Cramer-matched spectral correction, not independently introduced
Green fields. They imply
\[
\mathcal L_sw=F-\delta u_0,\qquad
(\mathcal L_s+\delta)w=F-\delta u_1.
\]
Differentiating \emph{along} \(r=\rho(s)\), followed by use of
the original equation at the endpoint, gives
\begin{align*}
(u_1)_s(s,\rho)&=0,&
(u_1)_{sr}(s,\rho)&=-\rho'\delta U,\\
w_s(s,\rho)&=U'-\rho'D,&
w_{sr}(s,\rho)&=
 D'-\rho'\{F(s,\rho)-p(s,\rho)D-(q(s,\rho)+\delta)U\}.
\end{align*}
In particular the original regularized derivative jet obeys
\[
\rho^2(u_1)_{sr}(s,\rho)+\rho(u_1)_s(s,\rho)
=-\rho'\rho^2\delta U .
\]
For the corrected physical field the full identity is
\[
\begin{split}
\rho^2w_{sr}(s,\rho)+\rho w_s(s,\rho)
={}&\rho^2D'+\rho U'\\
&-\rho'
\{\rho^2F-\rho^2pD-\rho^2(q+\delta)U+\rho D\}_{r=\rho}.
\end{split}
\]
For the actual moving regularized Kummer operator
\(p=2/r+i\alpha(s)r/2\), the braces simplify to
\(\rho^2F-(\rho+i\alpha\rho^3/2)D-\rho^2(q+\delta)U\).
No explicit inverse power of the shrinking left radius is introduced.

At a \emph{fixed} Kummer radius, direct operator
differentiation, without independently differentiating the fundamental
basis, gives the triangular system
\[
\begin{aligned}
\mathcal L_s(u_0)_s
 &=F_s-p_s(u_0)_r-q_su_0,\\
\mathcal L_s(u_1)_s
 &=\delta'u_0+\delta(u_0)_s
 -p_s(u_1)_r-q_su_1.
\end{aligned}
\]
Let \(h_s\) be the original homogeneous Kummer initial-value mode
with \(h_s(\rho)=0\), \(\partial_rh_s(\rho)=1\).
Then \((u_1)_s+\rho'\delta U h_s\) has zero left value
\emph{and} slope, and solves the second inhomogeneous equation above.
Only this modified response is eligible for direct application of
the existing zero-left weighted Volterra estimate. The homogeneous
mode's norm, the actual source derivative, and the selected
spectral/compatibility derivative remain to be established.

\noindent\textbf{Necessity of a genuine temporal producer.}
For a fixed nonzero original Green response
\(h=\mathcal G_\rho[u_0]\), choose
\(\delta_N(s)=\varepsilon\sin(Ns)\).
Then \(\sup_s\|\delta_N(s)h\|_X
\le\varepsilon\|h\|_X\), with identical
zero-left \(C^1\) spatial jets, while
\(\|\partial_s(\delta_Nh)(0)\|_X
=\varepsilon N\|h\|_X\).
Thus the existing frozen-time theorem types alone cannot imply
a time-positive residual. The \texttt{delta : Real -> Complex}
field of the actual Lean moving packet has no derivative constraint.
This example excludes an invalid \emph{logical inference}; it does
not exclude a physically selected Hartree trajectory.

\noindent\textbf{Verification boundary.}
This is a paper calculation only and has no new Lean acceptance.
CI 43124289518 strictly verifies the original four-region
nonlinear PDE telescoping and three spatial \(C^1\) joins, not
this time-differentiated Green system. The same-orbit
\(F_s,\delta',\rho',p_s,q_s\) estimates, full untruncated Hartree
source, parameter Lipschitz, real finite recursion and
PDE-derived Type--II rate remain open. The retained original
environment count is \(11/13\).
